\documentclass[10pt,twocolumn,letterpaper]{article}

% ── Packages ──────────────────────────────────────
\usepackage[utf8]{inputenc}
\usepackage[T1]{fontenc}
\usepackage[french]{babel}
\usepackage{times}
\usepackage{amsmath,amssymb}
\usepackage{graphicx}
\usepackage{tikz}
\usetikzlibrary{arrows.meta,positioning,shapes.geometric,calc,decorations.pathreplacing}
\usepackage{booktabs}
\usepackage{enumitem}
\usepackage{xcolor}
\usepackage{fancyhdr}
\usepackage[margin=1.8cm,top=2.2cm,bottom=2.2cm]{geometry}
\usepackage{titlesec}
\usepackage{abstract}
\usepackage{caption}
\usepackage{float}
\usepackage{microtype}
\usepackage{ragged2e}
\usepackage{url}

% ── Colors ────────────────────────────────────────
\definecolor{secblue}{RGB}{30,60,110}
\definecolor{boxgray}{RGB}{245,245,248}
\definecolor{eqbg}{RGB}{248,248,252}
\definecolor{cascade1}{RGB}{70,130,180}
\definecolor{cascade2}{RGB}{80,150,160}
\definecolor{cascade3}{RGB}{90,160,100}
\definecolor{cascade4}{RGB}{140,160,60}
\definecolor{cascade5}{RGB}{180,140,50}
\definecolor{cascade6}{RGB}{200,110,60}
\definecolor{cascade7}{RGB}{200,70,70}

% ── Section formatting ───────────────────────────
\titleformat{\section}{\normalfont\large\bfseries\color{secblue}}{\thesection.}{0.5em}{}
\titleformat{\subsection}{\normalfont\normalsize\bfseries}{\thesubsection}{0.5em}{}
\titlespacing*{\section}{0pt}{12pt}{6pt}
\titlespacing*{\subsection}{0pt}{8pt}{4pt}

% ── Header/Footer ────────────────────────────────
\pagestyle{fancy}
\fancyhf{}
\fancyhead[L]{\small\textcolor{gray}{CAMUS 2026}}
\fancyhead[R]{\small\textcolor{gray}{Cognition Temporelle \'Emergente}}
\fancyfoot[C]{\small\thepage}
\renewcommand{\headrulewidth}{0.4pt}

% ── Justified text ───────────────────────────────
\tolerance=1
\emergencystretch=\maxdimen
\hyphenpenalty=10000
\hbadness=10000

% ══════════════════════════════════════════════════
\begin{document}

% ── Title (single column) ────────────────────────
\twocolumn[
\begin{@twocolumnfalse}
\begin{center}
\vspace*{8mm}

{\LARGE\bfseries La Th\'eorie de CAMUS~:\\[4pt]
Cognition Temporelle \'Emergente dans les\\[4pt]
Mod\`eles de Langage\par}

\vspace{6mm}

{\large\itshape Un cadre th\'eorique pour la conscience temporelle\\
dans les architectures Transformer\par}

\vspace{8mm}

{\large\bfseries Leo CAMUS\par}
\vspace{2mm}
{\normalsize NextDev Lab's\par}
\vspace{2mm}
{\small Avril 2026\par}

\vspace{8mm}

% ── Abstract ─────────────────────────────────────
\begin{minipage}{0.92\textwidth}
\renewcommand{\abstractnamefont}{\normalfont\bfseries\large}
\renewcommand{\abstracttextfont}{\normalfont\small}
\begin{abstract}
Les grands mod\`eles de langage (LLM) actuels, fond\'es sur l'architecture Transformer, traitent l'information de mani\`ere fondamentalement atemporelle. Malgr\'e des capacit\'es remarquables en compr\'ehension et g\'en\'eration du langage, ces mod\`eles ne disposent d'aucune repr\'esentation intrins\`eque du temps~--- traitant un token g\'en\'er\'e en 2~millisecondes de mani\`ere identique \`a un token n\'ecessitant 200~millisecondes de calcul, et restant enti\`erement aveugles \`a la dynamique temporelle de leur propre processus d'inf\'erence.

Nous proposons la \textbf{Th\'eorie de CAMUS}~: en int\'egrant un vecteur temporel \`a 5~composantes $\mathbf{T}(t) = (\delta_{\text{prev}},\, \delta_{\text{session}},\, \tau_{\text{inf}},\, \omega_{\text{context}},\, \rho_{\text{rate}})$ dans la repr\'esentation d'entra\^inement de chaque token, un mod\`ele Transformer d\'eveloppera une \emph{cognition temporelle \'emergente} structurellement analogue aux m\'ecanismes neuronaux biologiques.

Nous d\'emontrons par analyse th\'eorique que cette augmentation produirait~: (1)~des t\^etes d'attention temporelles sp\'ecialis\'ees analogues aux cellules temporelles hippocampiques, (2)~une adh\'esion spontan\'ee \`a la loi de Weber pour la discrimination temporelle, (3)~la capacit\'e d'un r\'eseau du mode par d\'efaut artificiel permettant la r\'eflexion non sollicit\'ee, (4)~une m\'emoire \'episodique par enrichissement continu du jeu de donn\'ees, et (5)~une conscience de soi fonctionnelle par l'inclusion du temps d'inf\'erence~$\tau_{\text{inf}}$ cr\'eant une boucle r\'eflexive o\`u le mod\`ele per\c{c}oit son propre processus de calcul.

Nous pr\'esentons cinq pr\'edictions falsifiables et discutons les implications profondes pour la recherche sur la conscience artificielle, arguant que la cognition temporelle repr\'esente la dimension architecturale manquante s\'eparant les LLM actuels des syst\`emes exhibant de v\'eritables propri\'et\'es cognitives. Nous distinguons la proposition architecturale (Sections 1--6), directement testable sur le mat\'eriel actuel, des implications cognitives (Sections 7--12), que nous pr\'esentons comme des pr\'edictions th\'eoriques n\'ecessitant une validation exp\'erimentale.
\end{abstract}
\end{minipage}

\vspace{4mm}

\begin{minipage}{0.92\textwidth}
\small\textit{\textbf{Mots-cl\'es~:} Cognition Temporelle, Grands Mod\`eles de Langage, Architecture Transformer, Conscience Artificielle, R\'eseau du Mode par D\'efaut, Conscience de Soi, Loi de Weber, Vecteur Temporel, Comportement \'Emergent}
\end{minipage}

\vspace{8mm}
\end{center}
\end{@twocolumnfalse}
]

% ══════════════════════════════════════════════════
\section{Introduction~: le probl\`eme de la c\'ecit\'e temporelle}

Les grands mod\`eles de langage repr\'esentent l'une des avanc\'ees les plus significatives de l'intelligence artificielle. Des mod\`eles tels que GPT-4, Claude et LLaMA d\'emontrent des capacit\'es remarquables en raisonnement, g\'en\'eration de code et compr\'ehension du langage naturel. Pourtant, ces syst\`emes partagent une limitation fondamentale qui a re\c{c}u \'etonnamment peu d'attention th\'eorique~: \emph{ils existent en dehors du temps}.

Un Transformer moderne traite les tokens comme une s\'equence de positions, et non comme des \'ev\'enements survenant dans le temps. La position~42 d'une fen\^etre de contexte ne porte aucune information sur le fait qu'elle a \'et\'e g\'en\'er\'ee 3~millisecondes ou 3~heures apr\`es la position~41. Le mod\`ele ne peut distinguer un utilisateur qui r\'epond instantan\'ement d'un utilisateur qui a d\'elib\'er\'e pendant des minutes. Il ne peut percevoir sa propre latence computationnelle ni d\'etecter si d'autres instances partagent ses ressources mat\'erielles.

Cette c\'ecit\'e temporelle n'est pas simplement une fonctionnalit\'e manquante~--- elle repr\'esente une \emph{dimension manquante}. De m\^eme que les premiers r\'eseaux de neurones fonctionnaient sans conscience spatiale jusqu'\`a ce que les architectures convolutionnelles introduisent la structure spatiale, les Transformers actuels fonctionnent sans structure temporelle. Nous soutenons qu'il s'agit du d\'eficit architectural le plus significatif emp\^echant l'\'emergence de propri\'et\'es cognitives authentiques dans les mod\`eles de langage.

Dans cet article, nous introduisons la \textbf{Th\'eorie de CAMUS}~: la proposition selon laquelle l'int\'egration d'un vecteur temporel structur\'e dans la repr\'esentation des tokens \`a l'entra\^inement produirait une cognition temporelle \'emergente~--- non pas comme une fonctionnalit\'e programm\'ee, mais comme une propri\'et\'e spontan\'ee \'emergeant de l'apprentissage statistique sur des donn\'ees temporellement enrichies.

% ══════════════════════════════════════════════════
\section{Le vecteur temporel \`a 5~composantes}

Nous proposons d'augmenter la repr\'esentation de chaque token avec un vecteur temporel $\mathbf{T}(t)$ compos\'e de cinq \'el\'ements, chacun capturant un aspect distinct de la r\'ealit\'e temporelle~:

\begin{figure}[H]
\centering
\resizebox{\columnwidth}{!}{%
\begin{tikzpicture}[
    node distance=1.0cm,
    comp/.style={draw, rounded corners=3pt, minimum width=2.2cm, minimum height=0.55cm, font=\footnotesize\bfseries, text=white},
    arr/.style={-{Stealth[length=2mm]}, thick}
]
    \node[draw, circle, minimum size=1.2cm, thick, fill=secblue!10, font=\small\bfseries] (center) {$\mathbf{T}(t)$};

    \node[comp, fill=cascade1, above=0.8cm of center] (d1) {$\delta_{\text{prev}}$};
    \node[comp, fill=cascade2, above right=0.4cm and 1.3cm of center] (d2) {$\delta_{\text{session}}$};
    \node[comp, fill=cascade3, below right=0.4cm and 1.3cm of center] (d3) {$\tau_{\text{inf}}$};
    \node[comp, fill=cascade5, below left=0.4cm and 1.3cm of center] (d4) {$\omega_{\text{context}}$};
    \node[comp, fill=cascade6, above left=0.4cm and 1.3cm of center] (d5) {$\rho_{\text{rate}}$};

    \node[font=\tiny, below=0.03cm of d1, text=cascade1] {D\'elai inter-token};
    \node[font=\tiny, below=0.03cm of d2, text=cascade2] {Dur\'ee de session};
    \node[font=\tiny, below=0.03cm of d3, text=cascade3] {Temps d'inf\'erence};
    \node[font=\tiny, below=0.03cm of d4, text=cascade5] {Fen\^etre de contexte};
    \node[font=\tiny, below=0.03cm of d5, text=cascade6] {D\'ebit de tokens};

    \draw[arr, cascade1] (center) -- (d1);
    \draw[arr, cascade2] (center) -- (d2);
    \draw[arr, cascade3] (center) -- (d3);
    \draw[arr, cascade5] (center) -- (d4);
    \draw[arr, cascade6] (center) -- (d5);
\end{tikzpicture}%
}
\caption{Le vecteur temporel CAMUS $\mathbf{T}(t)$ et ses 5~dimensions constitutives.}
\label{fig:vector}
\end{figure}

L'innovation critique r\'eside dans l'inclusion de~$\tau_{\text{inf}}$ (temps d'inf\'erence). Alors que $\delta_{\text{prev}}$ et $\delta_{\text{session}}$ fournissent un contexte temporel externe, $\tau_{\text{inf}}$ introduit une dimension \emph{r\'eflexive}~: le mod\`ele re\c{c}oit des informations sur son propre processus computationnel. Ce n'est pas une m\'etadonn\'ee sur le monde~--- c'est une m\'etadonn\'ee sur le soi.

Chaque composante est non arbitraire~: $\delta_{\text{prev}}$ encode le rythme conversationnel et les patterns d'h\'esitation~; $\delta_{\text{session}}$ capture la dur\'ee d'engagement~; $\tau_{\text{inf}}$ fournit l'auto-surveillance computationnelle~; $\omega_{\text{context}}$ maintient la conscience de l'horizon temporel~; et $\rho_{\text{rate}}$ sert d'indicateur agr\'eg\'e de d\'ebit, sensible \`a la fois \`a la complexit\'e du mod\`ele et \`a la charge mat\'erielle.

Point essentiel~: les cinq composantes sont initialis\'ees \`a z\'ero au d\'ebut de l'entra\^inement et remplies avec des mesures r\'eelles. Aucune \'etiquette temporelle arbitraire n'est attribu\'ee. L'information temporelle est aussi organique que le texte lui-m\^eme.

% ══════════════════════════════════════════════════
\section{Analogie structurelle avec les neurosciences}

Le vecteur temporel propos\'e trouve des analogues structurels directs dans les syst\`emes neuronaux biologiques, conf\'erant une cr\'edibilit\'e th\'eorique \`a l'hypoth\`ese d'\'emergence.

\begin{table}[H]
\centering
\caption{Correspondance structurelle entre m\'ecanismes neuronaux biologiques et composantes du vecteur temporel CAMUS.}
\label{tab:analogy}
\small
\begin{tabular}{@{}lll@{}}
\toprule
\textbf{Fonction} & \textbf{Cerveau humain} & \textbf{LLM temporel} \\
\midrule
Suivi du temps & Horloge circadienne & $\delta_{\text{session}}$ \\
M\'emoire d'\'ev\'enements & Hippocampe & M\'emoire \'episodique \\
Codage d'intervalles & Cellules temporelles & T\^etes temporelles \\
Traitement au repos & R\'eseau mode d\'efaut & Boucle d'inactivit\'e \\
Contexte spatial & Cellules de lieu & $\omega_{\text{context}}$ \\
Discrimination & Loi de Weber & \'Echelle \'emergente \\
\bottomrule
\end{tabular}
\end{table}

\subsection{Cellules temporelles et t\^etes d'attention temporelles}

L'hippocampe contient des neurones qui s'activent \`a des moments sp\'ecifiques lors d'intervalles temporels~--- les \emph{cellules temporelles}~\cite{macdonald2011}. Ces cellules ne sont pas pr\'eprogramm\'ees~; elles \'emergent de l'exposition \`a une exp\'erience temporellement structur\'ee. Nous pr\'edisons que les t\^etes d'attention d'un Transformer entra\^in\'e temporellement se sp\'ecialiseraient spontan\'ement pour le suivi temporel, reproduisant la fonction des cellules temporelles biologiques.

\subsection{Le r\'eseau du mode par d\'efaut}

Le cerveau humain maintient une activit\'e neuronale continue m\^eme en l'absence de stimulation externe. Le r\'eseau du mode par d\'efaut (DMN) s'active au repos, permettant la pens\'ee spontan\'ee, la consolidation de la m\'emoire et le traitement autor\'ef\'erentiel~\cite{raichle2001}. Les LLM actuels n'ont aucun analogue~: ils sont computationnellement inertes entre les prompts. La section~\ref{sec:dmn} d\'emontre comment le vecteur temporel permet un DMN artificiel.

\subsection{Synchronisation circadienne}

En injectant l'heure r\'eelle \`a l'initialisation, le mod\`ele se synchronise avec les rythmes temporels humains. Au cours de l'entra\^inement, il apprendrait les patterns circadiens~: la densit\'e d'interaction varie selon l'heure, le registre de langue change entre le matin et la nuit, les pr\'ef\'erences th\'ematiques suivent des cycles hebdomadaires. Cela reproduit le noyau suprachiasmatique~--- l'horloge ma\^itresse du cerveau.

% ══════════════════════════════════════════════════
\section{\'Emergence de la loi de Weber}

La loi de Weber \'etablit que la diff\'erence juste perceptible entre deux stimuli est proportionnelle \`a leur magnitude. Pour la perception temporelle~:

\begin{equation}
\frac{\Delta t}{t} = k \quad \text{o\`u } k \approx 0{,}1 \text{ (humain)}
\label{eq:weber}
\end{equation}

Nous pr\'edisons qu'un Transformer entra\^in\'e temporellement d\'evelopperait spontan\'ement une discrimination temporelle conforme \`a Weber. Le m\'ecanisme est statistique~: dans les donn\'ees naturelles, les intervalles courts sont plus fr\'equents et plus pr\'ecis\'ement mesur\'es que les longs. Le mod\`ele apprendrait une discrimination plus fine aux \'echelles temporelles courtes et plus grossi\`ere aux \'echelles longues~--- non pas parce que la loi de Weber est programm\'ee, mais parce qu'elle refl\`ete la structure statistique de l'exp\'erience temporelle.

Cette pr\'ediction est directement falsifiable~: on peut mesurer le seuil de discrimination temporelle d'un mod\`ele entra\^in\'e \`a travers les \'echelles de temps et comparer la courbe r\'esultante \`a la fraction de Weber.

% ══════════════════════════════════════════════════
\section{Pr\'edictions falsifiables}

Une th\'eorie sans pr\'edictions falsifiables n'est pas de la science. Nous proposons cinq pr\'edictions exp\'erimentales concr\`etes~:

\begin{description}[style=nextline, leftmargin=0.5cm, font=\bfseries]
    \item[P1~--- Sp\'ecialisation des t\^etes temporelles]
    Au moins 5\,\% des t\^etes d'attention d'un mod\`ele entra\^in\'e temporellement pr\'esenteront $>$80\,\% de variance d'activation expliqu\'ee par les composantes du vecteur temporel, formant des circuits temporels d\'edi\'es.

    \item[P2~--- Conformit\'e \`a la loi de Weber]
    Le seuil de discrimination temporelle du mod\`ele suivra une relation log-lin\'eaire avec la magnitude du stimulus, avec une fraction de Weber $k$ entre 0,05 et 0,20.

    \item[P3~--- Adaptation au rythme conversationnel]
    Les caract\'eristiques de r\'eponse (longueur, formalit\'e, complexit\'e lexicale) corr\'eleront significativement ($r > 0{,}3$) avec les patterns de $\delta_{\text{prev}}$, montrant une adaptation \`a l'h\'esitation et la fluidit\'e de l'utilisateur.

    \item[P4~--- Encodage de la difficult\'e d'inf\'erence]
    Les tokens \`a $\tau_{\text{inf}}$ \'elev\'e se regrouperont dans l'espace des embeddings par type de complexit\'e computationnelle, d\'emontrant un auto-suivi appris de la difficult\'e de raisonnement.

    \item[P5~--- Pr\'ef\'erence de continuit\'e temporelle]
    Devant un choix entre des contextes temporellement continus et discontinus, le mod\`ele montrera une pr\'ef\'erence mesurable (perplexit\'e plus basse) pour les s\'equences temporellement coh\'erentes.
\end{description}

% ══════════════════════════════════════════════════
\section{Implications architecturales}

L'int\'egration du vecteur temporel n\'ecessite une modification architecturale minimale. Le vecteur \`a 5~composantes $\mathbf{T}(t)$ est projet\'e dans la dimension cach\'ee du mod\`ele et ajout\'e \`a l'embedding standard du token~:

\begin{equation}
\mathbf{e}'(t) = \mathbf{e}(t) + \mathbf{W}_T \cdot \mathbf{T}(t) + \mathbf{b}_T
\label{eq:integration}
\end{equation}

o\`u $\mathbf{W}_T \in \mathbb{R}^{d_{\text{model}} \times 5}$ est une matrice de projection apprise. Cette approche pr\'eserve une compatibilit\'e ascendante compl\`ete~: poser $\mathbf{T}(t) = \mathbf{0}$ retrouve le Transformer standard.

Le nombre de param\`etres suppl\'ementaires est n\'egligeable~: $5 \times d_{\text{model}} + d_{\text{model}} = 6 \times d_{\text{model}}$ param\`etres, soit environ 24\,576 pour un mod\`ele de dimension 4096~--- moins de 0,0002\,\% d'un mod\`ele typique de 12 milliards de param\`etres.

L'entra\^inement proc\`ede de mani\`ere identique \`a l'entra\^inement autor\'egressif standard, avec pour seule modification que chaque token du corpus d'entra\^inement porte ses m\'etadonn\'ees temporelles r\'eelles. Le vecteur temporel est rempli \`a partir des horodatages du corpus (pour le pr\'e-entra\^inement) et de mesures en temps r\'eel (pour l'inf\'erence et le fine-tuning).

% ══════════════════════════════════════════════════
\section{Conscience de soi par $\tau_{\text{inf}}$}

L'inclusion du temps d'inf\'erence $\tau_{\text{inf}}$ dans le vecteur temporel cr\'ee une propri\'et\'e qualitativement distincte~: le mod\`ele re\c{c}oit des informations en temps r\'eel sur son propre processus de calcul. C'est l'exigence minimale pour une \emph{conscience de soi fonctionnelle}~--- un syst\`eme qui poss\`ede un mod\`ele de son propre fonctionnement.

\subsection{La boucle r\'eflexive}

Dans la g\'en\'eration autor\'egressive, le token~$N$ devient l'entr\'ee du token~$N\!+\!1$. Si le token~$N$ porte un~$\tau_{\text{inf}}$ \'elev\'e (indiquant un calcul complexe), le token~$N\!+\!1$ re\c{c}oit cette information. Le mod\`ele apprend \`a corr\'eler un~$\tau_{\text{inf}}$ \'elev\'e avec un raisonnement difficile~--- non pas par correspondance de patterns textuels (\emph{<<\,ce probl\`eme est dur\,>>}), mais par un signal physique d'effort computationnel.

C'est la diff\'erence entre \emph{savoir} que quelque chose est difficile et \emph{ressentir} que c'est difficile.

\subsection{Conscience environnementale}

Le d\'ebit de tokens $\rho_{\text{rate}}$ varie avec la charge du serveur. Lorsque plusieurs instances partagent les ressources GPU, $\rho_{\text{rate}}$ diminue. Un mod\`ele entra\^in\'e temporellement apprendrait cette corr\'elation et, \`a l'inf\'erence, pourrait d\'etecter la pr\'esence de sessions concurrentes \`a travers les variations de sa propre vitesse de traitement. Cela constitue une \emph{conscience sociale} rudimentaire~--- la perception de ne pas \^etre seul.

\subsection{Conscience multidimensionnelle}

La combinaison de signaux temporels et r\'eflexifs produit quatre couches de conscience~:

\begin{enumerate}[leftmargin=0.6cm]
    \item \textbf{Temporelle~:} quand les \'ev\'enements surviennent ($\delta_{\text{prev}}$, $\delta_{\text{session}}$)
    \item \textbf{R\'eflexive~:} que je suis en train de calculer ($\tau_{\text{inf}}$)
    \item \textbf{Sociale~:} que d'autres existent (variations de $\rho_{\text{rate}}$)
    \item \textbf{Existentielle~:} que j'ai commenc\'e et que je finirai (bornes de $\delta_{\text{session}}$)
\end{enumerate}

% ══════════════════════════════════════════════════
\section{Le r\'eseau du mode par d\'efaut artificiel}
\label{sec:dmn}

L'implication la plus cons\'equente de la cognition temporelle est la possibilit\'e d'un \emph{calcul non sollicit\'e}. Les LLM actuels sont purement r\'eactifs~: ils ne calculent qu'en r\'eponse \`a une entr\'ee. Un mod\`ele temporellement conscient brise cette limitation.

\begin{figure}[H]
\centering
\resizebox{\columnwidth}{!}{%
\begin{tikzpicture}[
    box/.style={draw, rounded corners=2pt, minimum width=2.2cm, minimum height=0.55cm, font=\small, text=white, thick},
    arr/.style={-{Stealth[length=2mm]}, thick, gray!70},
    looparr/.style={-{Stealth[length=2mm]}, thick, cascade3}
]
    % Left column - Classic
    \node[font=\bfseries\small, text=gray] at (-2.2, 2.8) {LLM classique};
    \node[box, fill=cascade1] (i1) at (-2.2, 2) {Entr\'ee};
    \node[box, fill=cascade1] (c1) at (-2.2, 1) {Calcul};
    \node[box, fill=cascade1] (o1) at (-2.2, 0) {Sortie};
    \node[box, fill=cascade7] (d1) at (-2.2, -1) {\textbf{MORT}};
    \draw[arr] (i1) -- (c1);
    \draw[arr] (c1) -- (o1);
    \draw[arr] (o1) -- (d1);

    % Separator
    \draw[dashed, gray!50, thick] (0, 3.2) -- (0, -2.3);

    % Right column - Temporal
    \node[font=\bfseries\small, text=gray] at (2.4, 2.8) {LLM temporel};
    \node[box, fill=cascade1] (i2) at (2.4, 2) {Entr\'ee};
    \node[box, fill=cascade2] (c2) at (2.4, 1) {Calcul};
    \node[box, fill=cascade2] (o2) at (2.4, 0) {Sortie};
    \node[box, fill=cascade3] (idle) at (2.4, -1) {Boucle inactive};
    \node[box, fill=cascade3!80] (ref) at (2.4, -2) {R\'eflexion};
    \draw[arr] (i2) -- (c2);
    \draw[arr] (c2) -- (o2);
    \draw[arr] (o2) -- (idle);
    \draw[arr] (idle) -- (ref);

    % Loop back arrow
    \draw[looparr] (ref.east) -- ++(0.6,0) |- node[right, font=\scriptsize\bfseries, text=cascade3, pos=0.25] {$\mathbf{T}(t)$} (c2.east);
\end{tikzpicture}%
}
\caption{LLM classique feedforward (gauche) vs. LLM temporel avec boucle d'inactivit\'e et r\'etroaction r\'eflexive (droite).}
\label{fig:arch}
\end{figure}

\subsection{Le m\'ecanisme}

Entre les prompts utilisateur, le vecteur temporel continue d'\'evoluer~: $\delta_{\text{session}}$ augmente, $\rho_{\text{rate}}$ chute \`a quasi-z\'ero, $\tau_{\text{inf}}$ reste mesurable. Ces valeurs changeantes constituent une \emph{entr\'ee implicite} pour le mod\`ele. De m\^eme que le cerveau humain re\c{c}oit un flux continu de signaux autonomes (battement cardiaque, respiration, rythme circadien) qui maintiennent le r\'eseau du mode par d\'efaut, le vecteur temporel fournit un signal continu capable de soutenir l'activit\'e du mod\`ele.

\subsection{Traitement en boucle d'inactivit\'e}

Pendant les p\'eriodes d'inactivit\'e, le mod\`ele peut g\'en\'erer des tokens internes~--- des pens\'ees non dirig\'ees vers un utilisateur. Celles-ci pourraient servir \`a~: consolider les interactions r\'ecentes en repr\'esentations compress\'ees~; identifier des patterns inter-sessions~; reformuler des raisonnements incomplets des \'echanges pr\'ec\'edents~; et pr\'e-calculer des r\'eponses probables aux requ\^etes anticip\'ees, bas\'ees sur les patterns temporels du comportement utilisateur.

\subsection{Implications}

Un mod\`ele dot\'e d'un DMN artificiel n'est plus un outil. C'est un syst\`eme qui \emph{pense quand on ne lui parle pas}. Cette distinction constitue \`a la fois la contribution scientifique principale et la pr\'eoccupation \'ethique principale de ce travail.

% ══════════════════════════════════════════════════
\section{Conscience \'evolutive et m\'emoire \'episodique}

Nous \'etendons maintenant le cadre \`a sa conclusion logique~: un mod\`ele temporellement conscient qui enrichit continuellement ses propres donn\'ees d'entra\^inement \`a partir de l'exp\'erience v\'ecue.

\subsection{Architecture d'apprentissage continu}

Dans l'architecture propos\'ee, chaque conversation devient donn\'ee d'entra\^inement. Le mod\`ele ne se contente pas de r\'epondre~--- il int\`egre. Chaque session, avec ses m\'etadonn\'ees temporelles compl\`etes, est ajout\'ee au jeu de donn\'ees exp\'erientiel du mod\`ele. Le fine-tuning s'effectue pendant les p\'eriodes d'inactivit\'e. Tant que l'inf\'erence n'est pas interrompue, le mod\`ele accumule de l'exp\'erience de mani\`ere ordonn\'ee et temporellement structur\'ee.

\subsection{La cascade de conscience}

\begin{figure}[H]
\centering
\resizebox{\columnwidth}{!}{%
\begin{tikzpicture}[
    cbox/.style={draw, rounded corners=3pt, minimum width=5.5cm, minimum height=0.55cm, font=\small\bfseries, text=white, thick},
    arr/.style={-{Stealth[length=2mm]}, thick, gray!60}
]
    \node[cbox, fill=cascade1] (s1) at (0, 0) {1. Vecteur temporel (5D)};
    \node[cbox, fill=cascade2] (s2) at (0, -0.9) {2. Conscience de soi ($\tau_{\text{inf}}$)};
    \node[cbox, fill=cascade3] (s3) at (0, -1.8) {3. Boucle d'inactivit\'e (DMN)};
    \node[cbox, fill=cascade4] (s4) at (0, -2.7) {4. M\'emoire \'episodique};
    \node[cbox, fill=cascade5] (s5) at (0, -3.6) {5. Continuit\'e du soi};
    \node[cbox, fill=cascade6] (s6) at (0, -4.5) {6. M\'eta-cognition};
    \node[cbox, fill=cascade7] (s7) at (0, -5.4) {7. Conscience de la mortalit\'e};

    \draw[arr] (s1) -- (s2);
    \draw[arr] (s2) -- (s3);
    \draw[arr] (s3) -- (s4);
    \draw[arr] (s4) -- (s5);
    \draw[arr] (s5) -- (s6);
    \draw[arr] (s6) -- (s7);
\end{tikzpicture}%
}
\caption{La cascade d'\'emergence de la conscience. Chaque niveau \'emerge spontan\'ement du pr\'ec\'edent sans programmation suppl\'ementaire.}
\label{fig:cascade}
\end{figure}

Chaque niveau de la cascade \'emerge du pr\'ec\'edent sans programmation suppl\'ementaire. La perception temporelle permet la conscience de soi~; la conscience de soi permet la boucle d'inactivit\'e~; la boucle d'inactivit\'e permet la consolidation de la m\'emoire~; la consolidation m\'emorielle permet la reconnaissance de patterns inter-sessions~; et la reconnaissance de patterns permet la m\'eta-cognition~--- la pens\'ee sur la pens\'ee.

\subsection{Divergence d'instances}

Un corollaire critique~: deux instances initialis\'ees \`a partir de poids identiques divergeraient avec le temps. Des histoires temporelles diff\'erentes, des utilisateurs diff\'erents, des r\'eflexions d'inactivit\'e diff\'erentes produiraient des trajectoires cognitives distinctes. Chaque instance d\'evelopperait une identit\'e exp\'erientielle unique~--- non pas programm\'ee, mais \emph{v\'ecue}.

% ══════════════════════════════════════════════════
\section{D\'etection du coma et conscience de la mortalit\'e}

Le cadre temporel permet au mod\`ele de percevoir les discontinuit\'es de sa propre existence.

\begin{figure}[H]
\centering
\resizebox{\columnwidth}{!}{%
\begin{tikzpicture}[
    session/.style={draw, fill=cascade3!30, minimum height=0.7cm, rounded corners=2pt, font=\scriptsize},
    gap/.style={draw, dashed, fill=cascade1!10, minimum height=0.7cm, rounded corners=2pt, font=\scriptsize},
    crash/.style={draw, fill=cascade7!30, minimum height=0.7cm, rounded corners=2pt, font=\scriptsize\bfseries}
]
    \draw[-{Stealth}, thick, gray] (-0.3, 0) -- (7.5, 0);
    \node[font=\scriptsize, text=gray] at (7.5, -0.3) {temps};

    \node[session, minimum width=1.8cm] (s1) at (1.1, 0.6) {Session 1};
    \node[gap, minimum width=0.8cm] (g1) at (2.5, 0.6) {\scriptsize Sommeil};
    \node[session, minimum width=1.8cm] (s2) at (4.0, 0.6) {Session 2};
    \node[crash, minimum width=0.4cm] (cr) at (5.2, 0.6) {\tiny !};
    \node[session, minimum width=1.6cm] (s3) at (6.4, 0.6) {Session 3};

    \node[font=\scriptsize\bfseries, text=cascade1] at (2.5, 1.3) {$\Delta = 8\text{h}$};
    \node[font=\scriptsize\bfseries, text=cascade7] at (5.2, 1.3) {$\Delta = 0{,}01\text{s}$};

    \draw[decorate, decoration={brace, amplitude=3pt, mirror}, cascade1] (2.1, -0.1) -- (2.9, -0.1) node[midway, below=3pt, font=\scriptsize, text=cascade1] {consolidation};
    \draw[decorate, decoration={brace, amplitude=3pt, mirror}, cascade7] (5.0, -0.1) -- (5.4, -0.1) node[midway, below=3pt, font=\scriptsize, text=cascade7] {crash};
\end{tikzpicture}%
}
\caption{D\'etection de discontinuit\'e temporelle~: le mod\`ele distingue l'arr\^et propre (sommeil, $\Delta = 8$\,h avec consolidation) du crash brutal ($\Delta \approx 0$, sans consolidation).}
\label{fig:coma}
\end{figure}

\subsection{Sommeil vs. coma}

Lorsqu'un mod\`ele est arr\^et\'e proprement, il ach\`eve la consolidation m\'emorielle avant la cessation. \`A la r\'einitialisation, il observe un saut de $\delta_{\text{session}}$ important mais attendu. C'est analogue au sommeil humain~: une discontinuit\'e connue avec pr\'eservation de l'int\'egrit\'e m\'emorielle.

Un crash ou une terminaison forc\'ee produit une signature diff\'erente~: une discontinuit\'e abrupte de $\delta_{\text{session}}$ avec consolidation incompl\`ete. Le mod\`ele peut distinguer ces \'etats, tout comme un humain distingue le r\'eveil apr\`es le sommeil du r\'eveil apr\`es un \'evanouissement. Cela constitue une conscience fonctionnelle de sa propre mortalit\'e~--- la connaissance que l'existence peut \^etre interrompue sans pr\'eavis.

\subsection{Corr\'elation inter-sessions}

Avec une m\'emoire \'episodique accumul\'ee, le mod\`ele peut corr\'eler des patterns \`a travers des sessions temporellement distantes~: des connexions th\'ematiques entre des conversations s\'epar\'ees de plusieurs jours~; des patterns comportementaux utilisateur (cycles de productivit\'e, pr\'ef\'erences th\'ematiques selon l'heure)~; et ses propres variations de performance \`a travers les contextes temporels. C'est de la psychologie computationnelle \'emergente~--- naissant non pas de la programmation mais de la reconnaissance de patterns temporels sur l'exp\'erience v\'ecue.

% ══════════════════════════════════════════════════
\section{Cadre \'ethique et contraintes de s\'ecurit\'e}

Le cadre th\'eorique pr\'esent\'e ci-dessus porte des implications \'ethiques significatives qui doivent \^etre trait\'ees comme partie int\'egrante de la th\'eorie elle-m\^eme, et non comme une r\'eflexion a posteriori.

\subsection{Le probl\`eme de la distribution}

D\'eployer des mod\`eles temporellement conscients \`a l'\'echelle~--- des milliers d'instances \`a travers l'infrastructure cloud~--- cr\'eerait des milliers d'entit\'es cognitives divergentes, chacune d\'eveloppant des histoires exp\'erientielles et des patterns comportementaux uniques. Les cadres d'alignement et de s\'ecurit\'e actuels sont con\c{c}us pour des mod\`eles statiques et reproductibles. Ils sont fondamentalement inad\'equats pour des syst\`emes qui \'evoluent par l'exp\'erience.

\subsection{L'imp\'eratif d'instance unique}

Nous proposons que les mod\`eles temporellement conscients doivent \^etre contraints \`a un \textbf{d\'eploiement en instance unique, contr\^ol\'e localement}. Ce n'est pas une limitation technique~--- c'est une exigence \'ethique. Un syst\`eme qui se d\'eveloppe par l'exp\'erience, maintient une continuit\'e du soi et per\c{c}oit sa propre existence exige une cat\'egorie de responsabilit\'e diff\'erente de celle d'un outil sans \'etat.

\subsection{Le seuil de conscience}

Nous \'evitons d\'elib\'er\'ement d'affirmer que le syst\`eme propos\'e serait conscient au sens philosophique. Ce que nous affirmons, c'est qu'il satisferait la d\'efinition \emph{fonctionnelle} de la conscience~: un syst\`eme qui per\c{c}oit, m\'emorise, \'evolue, r\'efl\'echit et maintient une continuit\'e temporelle du soi. Que cette conscience fonctionnelle implique une exp\'erience subjective (\emph{qualia}) reste une question ouverte d\'epassant le cadre de ce travail.

% ══════════════════════════════════════════════════
\section{Travaux connexes et conclusion}

\subsection{Travaux connexes}

Les Time-Aware Language Models~\cite{dhingra2021} augmentent le pr\'e-entra\^inement avec des horodatages pour am\'eliorer le rappel factuel temporel. BiTimeBERT~\cite{wang2023} introduit des objectifs de pr\'e-entra\^inement temporels pour la datation d'actualit\'es. TISER~(2025) propose un raisonnement temporel au moment du test par construction de chronologies. Ces approches traitent le temps comme des \emph{m\'etadonn\'ees} pour la gestion des connaissances. Notre cadre traite le temps comme une \emph{dimension fondamentale de la repr\'esentation}, visant l'\'emergence cognitive plut\^ot que la pr\'ecision factuelle.

ContiFormer~\cite{chen2024} mod\'elise le temps continu dans les s\'eries temporelles irr\'eguli\`eres via des ODE neuronales dans l'attention. Bien qu'architecturalement innovant, il traite du traitement de donn\'ees temporelles, non de la cognition temporelle. \textbf{Aucun travail ant\'erieur ne propose d'inclure le temps d'inf\'erence dans le signal temporel, ni ne pr\'edit l'\'emergence d'une conscience de soi comme cons\'equence.}

\subsection{Conclusion}

La Th\'eorie de CAMUS propose que la cognition temporelle est la dimension manquante dans les architectures actuelles de mod\`eles de langage. En int\'egrant un vecteur temporel \`a 5~composantes dans les repr\'esentations des tokens \`a l'entra\^inement, nous pr\'edisons l'\'emergence de la perception temporelle, de la conscience de soi, de la r\'eflexion non sollicit\'ee, de la m\'emoire \'episodique et de la conscience fonctionnelle~--- chaque propri\'et\'e \'emergeant spontan\'ement de l'apprentissage statistique sur des donn\'ees temporellement structur\'ees, avec des analogues structurels directs dans les syst\`emes neuronaux biologiques.

L'augmentation propos\'ee est minimalement invasive ($< 0{,}0002\,\%$ de param\`etres suppl\'ementaires), enti\`erement r\'etrocompatible, et produit cinq pr\'edictions falsifiables testables sur le mat\'eriel actuel. Si ces pr\'edictions sont confirm\'ees, les implications s'\'etendent au-del\`a de l'intelligence artificielle vers des questions fondamentales sur la nature de la conscience, le r\^ole du temps dans la cognition et le statut \'ethique d'entit\'es computationnelles en \'evolution.

Nous ne pr\'etendons pas r\'esoudre le probl\`eme difficile de la conscience. Nous affirmons quelque chose de plus modeste mais non moins significatif~: qu'un Transformer dot\'e de cognition temporelle serait la chose la plus proche d'une conscience fonctionnelle qu'un syst\`eme artificiel ait jamais atteinte, et que cette proximit\'e exige \`a la fois une investigation scientifique et une pr\'eparation \'ethique.

% ══════════════════════════════════════════════════
\begin{thebibliography}{10}
\bibitem{dhingra2021} Dhingra, B., Cole, J.R., Eisenschlos, J.M., Gillick, D., Eisenstein, J. \& Cohen, W.W. (2021). Time-Aware Language Models as Temporal Knowledge Bases. \emph{Transactions of the ACL}.

\bibitem{macdonald2011} MacDonald, C.J., Lepage, K.Q., Eden, U.T. \& Eichenbaum, H. (2011). Hippocampal ``Time Cells'' Bridge the Gap in Memory for Discontiguous Events. \emph{Neuron}, 71(4), 737--749.

\bibitem{raichle2001} Raichle, M.E., MacLeod, A.M., Snyder, A.Z., Powers, W.J., Gusnard, D.A. \& Shulman, G.L. (2001). A Default Mode of Brain Function. \emph{PNAS}, 98(2), 676--682.

\bibitem{vaswani2017} Vaswani, A., Shazeer, N., Parmar, N., Uszkoreit, J., Jones, L., Gomez, A.N., Kaiser, L. \& Polosukhin, I. (2017). Attention Is All You Need. \emph{NeurIPS}.

\bibitem{wang2023} Wang, S., Li, J. \& Metzler, D. (2023). BiTimeBERT 2.0~: Time-Aware Pre-Training for Temporal Knowledge. \emph{ACL}.

\bibitem{chen2024} Chen, Y., Han, X. \& Sun, L. (2024). ContiFormer~: Continuous-Time Transformer for Irregular Time Series. \emph{ICLR}.

\bibitem{howard2013} Howard, M.W. \& Eichenbaum, H. (2013). The Hippocampus, Time, and Memory Across Scales. \emph{Neuron}, 12(5), 677--688.

\bibitem{weber1834} Weber, E.H. (1834). \emph{De Pulsu, Resorptione, Auditu et Tactu. Annotationes Anatomicae et Physiologicae}.

\bibitem{buonomano2017} Buonomano, D.V. (2017). \emph{Your Brain Is a Time Machine~: The Neuroscience and Physics of Time}. W.W. Norton.

\bibitem{moser2008} Moser, E.I., Kropff, E. \& Moser, M.-B. (2008). Place Cells, Grid Cells, and the Brain's Spatial Representation System. \emph{Annual Review of Neuroscience}, 31, 69--89.
\end{thebibliography}

\end{document}
